\section*{ЗАКЛЮЧЕНИЕ}
\addcontentsline{toc}{section}{ЗАКЛЮЧЕНИЕ}

Криптографические методы защиты информации являются фундаментом современной цифровой безопасности, обеспечивая конфиденциальность данных при их хранении и передаче. Стандарт симметричного шифрования AES (Advanced Encryption Standard), благодаря своей надежности и эффективности, остается основным инструментом защиты информации в коммерческих и государственных системах.

В рамках данной работы была спроектирована и реализована программная библиотека для шифрования и дешифрования данных по алгоритму AES-128 на языке программирования Common Lisp. Использование функционального подхода позволило создать модульный и структурированный код.

Основные результаты работы:

\begin{enumerate}
	\item Проведен анализ предметной области. Изучены теоретические основы симметричного блочного шифрования, математический аппарат конечных полей Галуа и спецификация стандарта AES.
	\item Разработана архитектура системы. Спроектирована модульная структура библиотеки и выбраны эффективные структуры данных (векторы) для хранения ключей и состояния шифра.
	\item Осуществлена программная реализация алгоритма. Разработаны функции расширения ключа (key-expansion), раундовых преобразований шифрования и дешифрования, а также механизм дополнения данных по стандарту PKCS\#7.
	\item Проведено тестирование разработанного программного обеспечения. Выполнена верификация работы библиотеки на наборах данных различной длины, подтвердившая корректность процессов шифрования и полного восстановления исходной информации.
\end{enumerate}

Все требования, объявленные в техническом задании, были полностью реализованы, все задачи, поставленные в начале разработки проекта, решены в полном объеме.

Готовый рабочий проект представляет собой функционирующую программную библиотеку на языке Common Lisp, готовую к использованию для защиты данных произвольной длины.